\chapter{未分类图片转写待审核稿}

\section*{1.5 1.6}

\begin{note}[1.5 1.6]

文本: 36不看了

\end{note}

\begin{note}[来源路径]

近世代数/近世代数/无标题/1 5 1 6 1ead2efa3f42800d8ce5e3bfbf4c3c5e.md

\end{note}

\section*{maki完美算术}

\begin{note}[来源上级主题]

教程 (\%E6\%95\%99\%E7\%A8\%8B\%201a0d2efa3f428048bfcff3fa4cfa7a94.md)

\end{note}

\begin{note}[maki完美算术]

待根据图片进一步人工校对。

\end{note}

\begin{note}[来源路径]

近世代数/近世代数/无标题/maki完美算术 1c4d2efa3f42808f9aa7ffb283e08616.md

\end{note}

\section*{命题2.10}

\begin{proposition}[命题2.10]

fir jt 2.10

HEN) 4, MRRMRRES, AAEDAR EAST AN-MER (MRK EH, BRK” TRE).

TRAPPE SERRA AH. ARE BH, CMNNHRREMKS HR. RNRAMBR

eA “BO? HE, BY R(A) C (A), R(A)RC (A). REM HIE, BATE —-PMLGKA. Rikr eR, ae (A),

UUFERRRPHBRI, RANMA GEL. Kracl, RHF ERRPHBRI MARU, Albrae (A).

MIE YT (A) ZRATH.

NTA RPC a eR, ABARIVMRARAM, e-toc RE TA. — ib, IB

WENA AIRE PICK a1.-++ dn ER, BAHARA HICK ME ETA. FAD, ASR ee fal LY,

RTA EY ea TACHI PER

\end{proposition}

\begin{note}[图片 OCR 摘录，待人工复核]
以下内容来自源图片识别，便于审核时对照：

\textbf{图片 1}（image.png）
\begin{Verbatim}[fontsize=\small]
fir jt 2.10
HEN) 4, MRRMRRES, AAEDAR EAST AN-MER (MRK EH, BRK” TRE).

TRAPPE SERRA AH. ARE BH, CMNNHRREMKS HR. RNRAMBR
eA “BO? HE, BY R(A) C (A), R(A)RC (A). REM HIE, BATE —-PMLGKA. Rikr eR, ae (A),
UUFERRRPHBRI, RANMA GEL. Kracl, RHF ERRPHBRI MARU, Albrae (A).
MIE YT (A) ZRATH.

NTA RPC a eR, ABARIVMRARAM, e-toc RE TA. — ib, IB
WENA AIRE PICK a1.-++ dn ER, BAHARA HICK ME ETA. FAD, ASR ee fal LY,
RTA EY ea TACHI PER
\end{Verbatim}

\end{note}

\begin{note}[来源路径]

近世代数/近世代数/无标题/命题2 10 209d2efa3f4280f7a74ac9f1898d47da.md

\end{note}

\section*{基本概念}

\begin{note}[基本概念]

子级 项目: 映射与集合基础知识 (\%E6\%98\%A0\%E5\%B0\%84\%E4\%B8\%8E\%E9\%9B\%86\%E5\%90\%88\%E5\%9F\%BA\%E7\%A1\%80\%E7\%9F\%A5\%E8\%AF\%86\%201aed2efa3f42800c8261da73ac826269.md), 代数运算 (\%E4\%BB\%A3\%E6\%95\%B0\%E8\%BF\%90\%E7\%AE\%97\%201aed2efa3f4280d597defe353a625299.md), 分类与关系 (\%E5\%88\%86\%E7\%B1\%BB\%E4\%B8\%8E\%E5\%85\%B3\%E7\%B3\%BB\%201aed2efa3f42802491cdeb2b56b5f1f7.md)

\end{note}

\begin{note}[来源路径]

近世代数/近世代数/无标题/基本概念 1c3d2efa3f4280648bdcd081c8508c78.md

\end{note}

\section*{定理4}

\begin{note}[来源上级主题]

第三章 (\%E7\%AC\%AC\%E4\%B8\%89\%E7\%AB\%A0\%201ecd2efa3f4280c7a6b0ea0de416c09f.md)

\end{note}

\begin{theorem}[定理4]

文本: 78页

\end{theorem}

\begin{note}[图片 OCR 摘录，待人工复核]
以下内容来自源图片识别，便于审核时对照：

\textbf{图片 1}（image.png）
\begin{Verbatim}[fontsize=\small]
EB 4 eS EAR N-TSFH, SER BWR ORER PRAF S
Kira UTE MWR) SATS WAL. FKAS=S. MARE-TS
RWWHAR, MAS ZR WTF.
\end{Verbatim}

\end{note}

\begin{note}[来源路径]

近世代数/近世代数/无标题/定理4 1edd2efa3f4280a1a7f6f5906ff4407d.md

\end{note}

\section*{当时月影已不复}

\begin{note}[来源上级主题]

教程 (\%E6\%95\%99\%E7\%A8\%8B\%201a0d2efa3f428048bfcff3fa4cfa7a94.md)

\end{note}

\begin{note}[当时月影已不复]

待根据图片进一步人工校对。

\end{note}

\begin{note}[来源路径]

近世代数/近世代数/无标题/当时月影已不复 1c4d2efa3f428027a49ef0a8e75ed791.md

\end{note}

\section*{教程}

\begin{note}[教程]

子级 项目: 当时月影已不复 (\%E5\%BD\%93\%E6\%97\%B6\%E6\%9C\%88\%E5\%BD\%B1\%E5\%B7\%B2\%E4\%B8\%8D\%E5\%A4\%8D\%201c4d2efa3f428027a49ef0a8e75ed791.md), maki完美算术 (maki\%E5\%AE\%8C\%E7\%BE\%8E\%E7\%AE\%97\%E6\%9C\%AF\%201c4d2efa3f42808f9aa7ffb283e08616.md), 蔷薇科普朗克 (\%E8\%94\%B7\%E8\%96\%87\%E7\%A7\%91\%E6\%99\%AE\%E6\%9C\%97\%E5\%85\%8B\%201c4d2efa3f42806c99bec64a1ecd7a78.md)

无标题

\end{note}

\begin{note}[来源路径]

近世代数/近世代数/无标题/教程 1a0d2efa3f428048bfcff3fa4cfa7a94.md

\end{note}

\section*{无标题}

\begin{note}[无标题]

待根据图片进一步人工校对。

\end{note}

\begin{note}[来源路径]

近世代数/近世代数/无标题/无标题 1b5d2efa3f428047a266cc3033c3e6d4.md

\end{note}

\section*{无标题}

\begin{note}[无标题]

待根据图片进一步人工校对。

\end{note}

\begin{note}[来源路径]

近世代数/近世代数/无标题/无标题 20ad2efa3f4280639292c088afa9cce1.md

\end{note}

\section*{杂知识}

\begin{note}[杂知识]

非空子集构成子群的充要条件 a，b∈H ab逆∈H

群的定义 乘法封闭 有逆元 有结合律 有单位元

（S) 是S生成的子群是包含S的所有群的交也就是最小的包含S的子群当S是一个数的时候写作(4)

aH= \{a+x|x∈H\} 1H=4H 1和4只是不同代表元

整环中I的一个有逆元的元叫单位 b=εa ε为单位 b是a的相伴元

单位以及相伴元叫作a的平凡因子，其余a的因子叫a的真因子

不可约元只有平凡因子

主理想环的每一个理想都是主理想 主理想环一定是一个唯一分解环，不可约元可以生成极大理想

\end{note}

\begin{note}[关联图片]
以下图片已纳入本条整理，当前版本优先依据 Markdown 文字整理，图片细节可在审核时对照：

1. image.png\\

\end{note}

\begin{note}[来源路径]

近世代数/近世代数/无标题/杂知识 20ed2efa3f4280118b08e4bdc28f1b75.md

\end{note}

\section*{第三章}

\begin{note}[第三章]

子级 项目: 加群定义 (\%E5\%8A\%A0\%E7\%BE\%A4\%E5\%AE\%9A\%E4\%B9\%89\%201ecd2efa3f4280c5a839d3c295d23a6b.md), 环的定义 (\%E7\%8E\%AF\%E7\%9A\%84\%E5\%AE\%9A\%E4\%B9\%89\%201ecd2efa3f42801e8a2dfcc3ca0e824e.md), 交换环与单位元 (\%E4\%BA\%A4\%E6\%8D\%A2\%E7\%8E\%AF\%E4\%B8\%8E\%E5\%8D\%95\%E4\%BD\%8D\%E5\%85\%83\%201ecd2efa3f4280739601e8d3434ada2a.md), 环中逆元 (\%E7\%8E\%AF\%E4\%B8\%AD\%E9\%80\%86\%E5\%85\%83\%201ecd2efa3f428081aa09e3990df0dc13.md), 零因子 (\%E9\%9B\%B6\%E5\%9B\%A0\%E5\%AD\%90\%201ecd2efa3f42800fb79fed095ff2429e.md), 零因子与消去律互相证明 (\%E9\%9B\%B6\%E5\%9B\%A0\%E5\%AD\%90\%E4\%B8\%8E\%E6\%B6\%88\%E5\%8E\%BB\%E5\%BE\%8B\%E4\%BA\%92\%E7\%9B\%B8\%E8\%AF\%81\%E6\%98\%8E\%201ecd2efa3f42809fa11dc863d390a502.md), 整环定义 (\%E6\%95\%B4\%E7\%8E\%AF\%E5\%AE\%9A\%E4\%B9\%89\%201ecd2efa3f4280109d5bd2a376069406.md), 除环 (\%E9\%99\%A4\%E7\%8E\%AF\%201edd2efa3f428092b9f2f5fd34528e4a.md), 域就是交换的除环 (\%E5\%9F\%9F\%E5\%B0\%B1\%E6\%98\%AF\%E4\%BA\%A4\%E6\%8D\%A2\%E7\%9A\%84\%E9\%99\%A4\%E7\%8E\%AF\%201edd2efa3f428082b256f54c7b69c4d8.md), 定理1无零因子环所有不为0的元对加法阶都是一样的 (\%E5\%AE\%9A\%E7\%90\%861\%E6\%97\%A0\%E9\%9B\%B6\%E5\%9B\%A0\%E5\%AD\%90\%E7\%8E\%AF\%E6\%89\%80\%E6\%9C\%89\%E4\%B8\%8D\%E4\%B8\%BA0\%E7\%9A\%84\%E5\%85\%83\%E5\%AF\%B9\%E5\%8A\%A0\%E6\%B3\%95\%E9\%98\%B6\%E9\%83\%BD\%E6\%98\%AF\%E4\%B8\%80\%E6\%A0\%B7\%E7\%9A\%84\%201edd2efa3f428098a9d1d9282007041a.md), 特征的定义  对加法所有非零元素共同的阶 (\%E7\%89\%B9\%E5\%BE\%81\%E7\%9A\%84\%E5\%AE\%9A\%E4\%B9\%89\%20\%E5\%AF\%B9\%E5\%8A\%A0\%E6\%B3\%95\%E6\%89\%80\%E6\%9C\%89\%E9\%9D\%9E\%E9\%9B\%B6\%E5\%85\%83\%E7\%B4\%A0\%E5\%85\%B1\%E5\%90\%8C\%E7\%9A\%84\%E9\%98\%B6\%201edd2efa3f4280ff9276db7f4d563977.md), 定理2 整环，除环，域特征只能是素数或者无限 (\%E5\%AE\%9A\%E7\%90\%862\%20\%E6\%95\%B4\%E7\%8E\%AF\%EF\%BC\%8C\%E9\%99\%A4\%E7\%8E\%AF\%EF\%BC\%8C\%E5\%9F\%9F\%E7\%89\%B9\%E5\%BE\%81\%E5\%8F\%AA\%E8\%83\%BD\%E6\%98\%AF\%E7\%B4\%A0\%E6\%95\%B0\%E6\%88\%96\%E8\%80\%85\%E6\%97\%A0\%E9\%99\%90\%201edd2efa3f42802dbf6fd904ddc8051a.md), 子环与子除环 (\%E5\%AD\%90\%E7\%8E\%AF\%E4\%B8\%8E\%E5\%AD\%90\%E9\%99\%A4\%E7\%8E\%AF\%201edd2efa3f42803d86f8f66dab16dc23.md), 定理2环同态传递零元负元 如果是交换环传递单位元 (\%E5\%AE\%9A\%E7\%90\%862\%E7\%8E\%AF\%E5\%90\%8C\%E6\%80\%81\%E4\%BC\%A0\%E9\%80\%92\%E9\%9B\%B6\%E5\%85\%83\%E8\%B4\%9F\%E5\%85\%83\%20\%E5\%A6\%82\%E6\%9E\%9C\%E6\%98\%AF\%E4\%BA\%A4\%E6\%8D\%A2\%E7\%8E\%AF\%E4\%BC\%A0\%E9\%80\%92\%E5\%8D\%95\%E4\%BD\%8D\%E5\%85\%83\%201edd2efa3f42805cb845f45123227986.md), 定理3同构传递整环除环域的性质 (\%E5\%AE\%9A\%E7\%90\%863\%E5\%90\%8C\%E6\%9E\%84\%E4\%BC\%A0\%E9\%80\%92\%E6\%95\%B4\%E7\%8E\%AF\%E9\%99\%A4\%E7\%8E\%AF\%E5\%9F\%9F\%E7\%9A\%84\%E6\%80\%A7\%E8\%B4\%A8\%201edd2efa3f4280ea855efe8275531d49.md), 定理4 (\%E5\%AE\%9A\%E7\%90\%864\%201edd2efa3f4280a1a7f6f5906ff4407d.md), 7理想 (7\%E7\%90\%86\%E6\%83\%B3\%201edd2efa3f4280a88bfce32f4c292dfb.md), 8剩余类环同态与理想 (8\%E5\%89\%A9\%E4\%BD\%99\%E7\%B1\%BB\%E7\%8E\%AF\%E5\%90\%8C\%E6\%80\%81\%E4\%B8\%8E\%E7\%90\%86\%E6\%83\%B3\%201edd2efa3f42801295a9f3bc6dd211c2.md), 9极大理想 (9\%E6\%9E\%81\%E5\%A4\%A7\%E7\%90\%86\%E6\%83\%B3\%201efd2efa3f42802fa98fcbcc6a6ed8b2.md), 10商域 93 没讲 (10\%E5\%95\%86\%E5\%9F\%9F\%2093\%20\%E6\%B2\%A1\%E8\%AE\%B2\%201efd2efa3f42806fa86bd0829ae38236.md)

文本: 63-92（90）

\end{note}

\begin{note}[来源路径]

近世代数/近世代数/无标题/第三章 1ecd2efa3f4280c7a6b0ea0de416c09f.md

\end{note}

\section*{考前必看}

\begin{note}[考前必看]

a|b b=ak右大 犹大

|a|a的阶 n次对称群Sn的阶是n!

上下看映射123456

253461 1-2-5-6 34不变就不写 不是位置换是数字换

凯莱定理 任意一个群都同变换群同构

有限群和置换群同构 因为置换群的前提是有限集的若干个一一变化

子环注意验证一下

群单位元是加法 环单位元是乘法乘他等于自身的

无零因子环对加法的阶都一样 特征只能是无穷和素数

剩余类加环n是素数就可以进化成域

R/I是一个域当且仅当I是极大理想

[Z:H]是商群的元素个数 只有H是正规子群的时候才是商群Na=aN

零因子别找漏除了本来就是阶的因子的还有其他与群中元素相乘可以等于0的

剩余类群中与群的阶不互素的因子都是零因子

矩阵行列式为0都是零因子因为不可逆

中心就是最大的正规子群

S是n次对称群，分清S，Z

零因子:没有零因子有消去律，没有零因子加有限=有逆元

有限：元素都有阶

集和有限，结合律成立，消去律成立乘法封闭是群

有限给出单位元 自己映射到自己产生一个变化 a是G中任意一个元素

通过消去律ax=ag得到x=g ax可以是集合中任意的数所以可以是e当他是e的时候就找到了逆元

环的加法单位元叫零元 乘法单位元才叫单位元

素数 阶是素数的群都是循环群 从群里面取一个元组成循环群 剩余类环+素数=整环无零因子环

无零因子环特征为无限要不是素数 素数生成的理想都是极大理想 两个互素的数可以与两个整数构成

nq+sp=1

理想：1关于加法封闭2 ar=ra交换律

主理想：(4)都是由它生成加法封闭有交换律 父环是交换环理想写成 ra+na n是整数

父环有单位元的时候(a)的元可以写成 xay求和的形式

交换环又有单位元直接写成ra

极大理想 1不等于父环 2不被其他理想包含 两步走I ≠R 包含I的只有R

剩余类环R/I 只有零理想和单位理想当且仅当I是极大理想

交换环+只有零理想和单位理想就是一个域

剩余类环R/I是域当且仅当I是极大理想 主理想的一个不可约元生成一个极大理想

素数阶的群都是循环群，n阶循环群都与Zn剩余类加群同构，剩余类加环+素数=域

2 指数是2的子群一定是不变子群 陪集不变子群本身 左陪集个数等于右陪集个数，陪集之间无交

一个群同他的商群同态。有限群大于2的阶是偶数，等于2的阶自己是自己的逆

G和G‘同态 核是G的一个正规子群是由映射到G’的单位元的原像构成 像imf是G‘的子群

满射是每一个像都有原像

循环群的子群还是循环群因为子群中最小的元素就能生成整个子群  循环群的商群还是循环群

有限群G ,H≤G |G|=|H|:(G:H)

\end{note}

\begin{note}[来源路径]

近世代数/近世代数/无标题/考前必看 20ed2efa3f42809dbc47ebd433ba4dcb.md

\end{note}

\section*{蔷薇科普朗克}

\begin{note}[来源上级主题]

教程 (\%E6\%95\%99\%E7\%A8\%8B\%201a0d2efa3f428048bfcff3fa4cfa7a94.md)

\end{note}

\begin{note}[蔷薇科普朗克]

待根据图片进一步人工校对。

\end{note}

\begin{note}[来源路径]

近世代数/近世代数/无标题/蔷薇科普朗克 1c4d2efa3f42806c99bec64a1ecd7a78.md

\end{note}

\section*{证明}

\begin{note}[来源上级主题]

命题2.1负元 0元的运算 (\%E5\%91\%BD\%E9\%A2\%982\%201\%E8\%B4\%9F\%E5\%85\%83\%200\%E5\%85\%83\%E7\%9A\%84\%E8\%BF\%90\%E7\%AE\%97\%201d8d2efa3f4280408c76c66618c847b3.md)

\end{note}

\begin{proof}

然后两边都减去a0证明a0=0 然后可以再证明0a=0

因为环的乘法对加法有分配律和结合律只要凑出0即可，证明两个相等就把他们移到一边再证明他是0

3就是把利用两次性质2

\end{proof}

\begin{note}[关联图片]
以下图片已纳入本条整理，当前版本优先依据 Markdown 文字整理，图片细节可在审核时对照：

1. image.png\\
2. image 1.png\\

\end{note}

\begin{note}[来源路径]

近世代数/近世代数/无标题/证明 1edd2efa3f4280e090e4d05710609507.md

\end{note}
